%------------------------------------------------------------------------------%
\section{Introdução}

\emph{Eratóstenes e a circunferência da Terra}

% \begin{flushright}
% \color{red}
% JOSÉ ROBERTO V. COSTA
% \end{flushright}

Eratóstenes nasceu em Cirene, Grécia, e morreu em Alexandria, Egito, no
terceiro século AEC. Ele era bibliotecário-chefe da famosa Biblioteca de
Alexandria, e foi lá que ele encontrou, num velho papiro, indicações de que ao
meio-dia de cada 21 de junho na cidade de Assuã (ou Syene, no grego antigo) 800
km ao sul de Alexandria, uma vareta fincada verticalmente no solo não produzia
sombra.

Cultura inútil, diriam alguns. Não para um homem observador como Eratóstenes.
Ele percebeu que o fenômeno não ocorria no mesmo dia e horário em Alexandria, e
pensou: Se o mundo é plano como uma mesa, então as sombras das varetas têm de ser
iguais. E se isto não acontece é porque a Terra deve ser curva!

%------------------------------------------------------------------------------%
\begin{figure}
  \centering
  \includegraphics[width=0.8\textwidth]{siena}
  \caption{Siena}
\end{figure}
%------------------------------------------------------------------------------%

Mais do que isso. Quanto mais curva fosse a superfície da Terra, maior seria a
diferença no comprimento das sombras.

O Sol deveria estar tão longe que seus raios de luz chegam à Terra paralelos.
Varetas fincadas verticalmente no chão em lugares diferentes lançariam sombras
de comprimentos distintos.

Eratóstenes decidiu fazer um experimento. Ele mediu o comprimento da sombra em
Alexandria ao meio-dia de 21 de junho, quando a vareta em Assuã, ao Sul do
Egito, não produzia sombra. Assim, ele obteve o ângulo $A$, conforme a figura
abaixo.

%------------------------------------------------------------------------------%
\begin{figure}
  \centering
  \includegraphics[height=4cm,width=0.4\textwidth]{siena2}
  \qquad
  \includegraphics[height=4cm,width=0.2\textwidth]{siena3}
  \caption{Siena 2 e 3}
\end{figure}
%------------------------------------------------------------------------------%

Preste atenção na figura acima. O ângulo $B$ terá o mesmo valor de
$A$, pois o desenho do experimento de Eratóstenes se reduz a uma
geometria muito simples:
\begin{quotation}
  \noindent
  se duas retas paralelas interceptam uma reta
  transversal, então os ângulos correspondentes são iguais.
\end{quotation}

As retas paralelas são os raios de luz do Sol e a reta transversal é a que
passa pelo centro da Terra e pela vareta em Alexandria. O ângulo $B$ (também
igual a \ang{7}), é a uma fração conhecida da circunferência da Terra e
corresponde à distância entre Assuã e Alexandria!

Eratóstenes sabia que essa distância valia cerca de \SI{800}{\kilo\meter}
e então pensou: \ang{7} são aproximadamente \(\sfrac{1}{50}\)
de uma circunferência (\ang{360}).
E isso corresponde a cerca de \SI{800}{\kilo\meter}.
Oitocentos quilômetros vezes cinquenta são quarenta mil quilômetros,
de modo que deve ser este o valor da circunferência da Terra.

Valor encontrado atualmente: cerca de \SI{40.072}{\kilo\meter}
ao longo da linha do equador.
Um erro muito pequeno para uma medida tão simples, e feita há tanto tempo! Com
a circunferência, podemos calcular o diâmetro e o raio ou ainda o volume e a
área da superfície, mediante fórmulas simples.

Repare que o conhecimento utilizado por Eratóstenes (retas paralelas cortadas
por uma transversal) é formalmente adquirido hoje nas aulas de geometria do
Ensino Fundamental.

% {\color{red}
%   Fonte: Costa, J. R. V. Eratóstenes e a circunferência da Terra. Astronomia no
%   Zênite, jul 2000. Disponível em:
%   \url{http://www.zenite.nu/eratostenes-e-a-circunferencia-da-terra}
% }

%------------------------------------------------------------------------------%